\begin{table}[!htbp]
\centering
\small
\caption{Excess-attention differences and temporal changes.}
\label{tab:memory-decay-key-results}
\renewcommand{\arraystretch}{1.15}
\begin{tabularx}{\linewidth}{@{}>{\raggedright\arraybackslash}p{0.175\linewidth}>{\raggedright\arraybackslash}p{0.15\linewidth}>{\raggedright\arraybackslash}p{0.155\linewidth}>{\centering\arraybackslash}p{0.145\linewidth}>{\raggedright\arraybackslash}X@{}}
\toprule
Result & Estimate & \makecell[l]{95\% CI /\\ uncertainty} & \makecell[b]{Median\\ decay-equivalent\\ years} & Interpretation \\
\midrule
Pooled film-linked excess premium & 19.12 excess-attention points & [18.86, 19.38] & -- & broad excess-attention difference \\
Film-embedded by snapshot (time-respecting) & 19.15 excess-attention points & [18.88, 19.41] & -- & time-respecting excess-attention difference \\
Prior-attention matched excess premium & 2.52 excess-attention points & [2.16, 2.87] two-way cluster bootstrap & 3.78 & matched observational comparison \\
Strict temporal event contrast & 2.45 excess-attention points & [0.62, 4.28] & 3.41 & conservative temporal check \\
Strict temporal post contrast & 2.98 excess-attention points & [1.17, 4.80] & 4.11 & conservative temporal check \\
Lower pre-film excess-attention contrast & 5.04 excess-attention points & [2.97, 6.87] integrated bootstrap & 6.88 [3.89, 10.22] & placebo-adjusted difference in change among songs in the lower pre-film excess-attention tercile \\
Exceeding the pre-film threshold & 18.0 percentage-point difference (39.5\% vs 21.5\%) & [6.7, 27.6] pp integrated bootstrap & -- & difference in above-threshold probability \\
\bottomrule
\end{tabularx}
\begin{flushleft}\footnotesize
Excess attention is observed Spotify popularity minus the fitted age/snapshot-specific expected attention from canonical collective-memory decay curves estimated on never-film age means. Intervals: HC1 for the broad and film-embedded rows; a two-way cluster bootstrap over treated and reused control song records for the matched residual, accounting for control reuse and repeated song identities; matched-pair standard errors for the strict temporal rows ($n=93$ pairs); integrated nested-bootstrap inference (song-level resampling, decay-curve refitting, pseudo-event reassignment, and rematching) for the lower-tercile excess-attention rows. Decay-equivalent years ask how many years younger a song would need to be, under the fitted expected-attention curve, to reach the shifted attention level; they are interpretive translations, not literal measures of memory recovered or randomized causal effects.
\end{flushleft}
\end{table}
